
@@M0@@ l'homomorphisme d'augmentation : @@M1@@.

\textbf{Proposition 6.1.} \uline{Soit} @@M2@@ \uline{un schéma quasi}-\uline{compact ayant un faisceau ample} @@M3@@. \uline{Alors} @@M4@@ \uline{est un nilidéal}.

Comme @@M5@@ est engendré par les éléments du type @@M6@@, où @@M7@@ est la classe d'un @@M8@@-Module localement libre de type fini @@M9@@, et @@M10@@ l'augmentation de @@M11@@, il suffit de montrer qu'un tel élément est nilpotent. La quasi-compacité de @@M12@@, et la commutation du @@M13@@ aux sommes directes finies permettent de se ramener au cas où @@M14@@ est de rang constant. Enfin, en passant au fibré de drapeaux de @@M15@@, ce qui donne une injection sur les @@M16@@ (3.1), et conserve l'hypothèse d'existence d'un faisceau ample, on peut supposer @@M17@@ inversible.

Soit @@M18@@ le dual de @@M19@@. Pour @@M20@@ assez grand, @@M21@@ est engendré par ses sections globales; on peut donc trouver un entier @@M22@@, et un homomorphisme surjectif : @@M23@@, donc (tensorisant par @@M24@@) un homomorphisme surjectif : @@M25@@. Le noyau de cet homomorphisme est localement libre de rang @@M26@@. On obtient donc, en notant @@M27@@ la classe de @@M28@@:
@@M29@@
Or @@M30@@. Par conséquent, pour tout faisceau inversible @@M31@@, l'élément @@M32@@ est nilpotent pour @@M33@@ assez grand, ou encore son produit @@M34@@ par @@M35@@ est nilpotent.

On peut alors écrire:
@@M36@@
et d'après ce qui précède, @@M37@@ et @@M38@@ sont nilpotents si on choisit @@M39@@ assez grand. Donc @@M40@@ est nilpotent.


